\clearpage
\section{Introduction}
\label{iii:sec:introduccion}

A constitutive relation organizes a system's response, but its
functional form and its coefficients raise different questions. Once
the coefficients are known, the equation permits responses to be
calculated; determining why those coefficients jointly have their
values requires describing the structure that fixes them. This work
addresses the second question through a mathematical construction of
two oriented channels. Its thesis is that the preceding structure
determines the equation and that evaluation preserves a recoverable
part of that structure. The electromagnetic vacuum provides a precise
realization of this thesis: permittivity, permeability, impedance and
speed are obtained as four related readouts of a single response.

The principal result comprises both direct construction and inversion.
On the contractive domain, the constitutive operator is positive,
its eigenvalues lie in a specified interval and its four coordinates
satisfy exact identities. Strict monotonicity of the response function
provides a unique inverse; with the sheets labelled, that inverse
recovers the channels and angular pair. This recovery composes three
descriptions: the vacuum response, the shape of the action ellipse
and the Barbero--Immirzi functional. The comparison therefore extends
beyond numerical values: it identifies the operations transmitting
orientation and spectral information between these descriptions.

\subsection{Generating kernel and provenance of the response}

The model is constructed from a formal mathematical core called
Triadic Modular Holography, henceforth HMT. Its first object,
\emph{All Possible Paths} (APP), combines paths and arithmetic
evaluations on the $81$ positions of the support
$X_9=(\mathbb Z/9\mathbb Z)^2$. Translations with generators
$\{(1,0),(-1,0),(0,1),(0,-1)\}$ define its additive Cayley graph,
a two-dimensional cyclic lattice with toroidal realization. Sum and
product evaluations act on this support; each preserves its residue
and quotient. The path structure and the evaluations thus perform
different functions within APP: displacements specify routes, while
the arithmetic sheets determine the data transported along them.

The second object, TRIT, is the oriented ternary signature
$\{-1,0,+1\}$, obtained through the balanced representation of
classes modulo $3$. Its quadratic realization distinguishes the
hyperbolic, parabolic and elliptic types, respectively. The third
object is the \emph{Topological Phase Kernel}, henceforth TPK:
it composes phase selection, cardinal transport and memory update.
The dynamics acts on positions, sheets, quotients, orientation,
prefixes and boundary data. Phase return preserves the accumulated
history, so a winding may close one coordinate without restoring
the complete state.

Nonadic prolongation transmits this structure through compatible
prefixes and boundaries. The coordinates $\pi_{\rm HMT}$,
$\varphi_{\rm HMT}$, $e_{\rm HMT}$ and $\alpha_{\rm HMT}$ are
generated within its joint development. The dodecaphasic record $K$
and its periodic evaluation enter the determination of $\alpha$ through dodecaphasic closure;
the complete analytic representation expresses that same coordinate
in the constructed tail class. The action sections and angular pair
are subsequently obtained. Their use as antecedents preserves their
proofs within the article. This sequence fixes the genealogy of the
channels before the constitutive response is evaluated.

Sheet orientation and areal--volumetric incidence perform complementary
roles. The former distinguishes route return from restoration of its
orientation; the latter arises from the mode calendar and its
incidence matrix. Chronological frequency, the measure of the symbolic
system and the regional boundary mark retain their own definitions.
These distinctions, developed in Section~\ref{iii:sec:hojas}, specify
the information reaching the angular transport and prevent its
replacement by an isolated scalar ratio.

\clearpage
\subsection{Constitutive determination and recovery of shape}

With $T=q_+P_++q_-P_-$ and $0<q_\pm<1$, the temporal incidences
$90$ and $120$ determine
\[
\mathcal V=(I-T^{90})(I-T^{120})^{-1}.
\]
Writing $S=T^{30}$, the equation
$\mathcal V(I+S+S^2+S^3)=I+S+S^2$ exhibits the rule fixing the
response for this transport. Section~\ref{iii:sec:constitutiva}
proves existence and uniqueness and establishes the scalar function
$f:(0,1)\to(3/4,1)$. Its values $r_\pm=f(q_\pm^{30})$ produce
\[
\widehat\varepsilon=r_+^2,\qquad
\widehat\mu=r_-^2,\qquad
\widehat Z=r_-/r_+,\qquad
\widehat c=(r_+r_-)^{-1}.
\]
Positivity and the product and quotient readouts determine the two
squares once the constitutive relations are fixed. The asserted
uniqueness belongs to this explicit system of rules and its domain.

Cubic inversion recovers $q_\pm$ and then the coordinates
$x=-\frac12\log(q_+q_-)$,
$y=\frac12\log(q_-/q_+)$. With sheet labels preserved, impedance
and speed suffice for this reconstruction.
Proposition~\ref{iii:prop:forma-LC} subsequently obtains the
anisotropy of the action ellipse and the relative
inductive--capacitive impedance. Action, frequency and impedance
references complete its dimensional realization.
Proposition~\ref{iii:prop:identidad-velocidad} proves
$c_{\rm int}=c_{\rm vac}$ in the same chart and relates path
length to elementary duration. Inversion thus reconstructs specified
spectral antecedents, while route and memory remain in the enriched
state from which they arise.

This inverse also induces a composition law between responses:
multiplication of transports is expressed by the transported operation
and admits an additive coordinate. Section~\ref{amp-higgs-seccion}
then composes the same channels with the angular gauge and Higgs
signatures. The joint inverse, its exact domain, and impedance
reconstruction retain the scale, refinement, and normalization
conditions of that realization.

\subsection{Cycle, constants and oriented functionals}

The function $f$ admits a realization on the twelve-phase cycle.
Its subspaces of dimensions $3$ and $4$ produce the determinant
quotient; their normalized trace difference gives $-1/12$.
On an oriented plane of the same cycle, a quarter-turn $J$ satisfies
$J^2=-I$ and its resolvent determines $s/(1+s^2)$. Two logarithmic
integrations, with their conditions at the origin, generate the moment
$\mathcal C_2(s)$ and its endpoint value, Catalan's constant.
Theorem~\ref{iii:thm:catalan-pell} relates evaluation at
$2-\sqrt3$ exactly to that endpoint value and
$\log(2+\sqrt3)$. The algebraic Pell section and the angular
arguments $q_\pm^{30}$ are distinct evaluations of the same collection.

The traces of the transport powers also determine a periodic table that
reconstructs the logarithm of the constitutive quotient. Its moment
series, Mellin representation, and remainder bounds are proved in
Section~\ref{amp-afin-sec}. Diagonal and tensorial composition make
explicit the depth retained by each operation; spectral return remains
distinct from return of the state.

Theorem~\ref{iii:thm:restitucion-funcional-catalan} completes the
relation through integral reconstruction: the constitutive function
determines the entire oriented moment, and the zero-limit condition
at the origin fixes its normalization. The proof distinguishes recovery
of that collection from inversion of its two channel evaluations.
Both operations preserve the speed readout
$\widehat c=(\det\mathcal V)^{-1}$ and its normalized trace, linking
constitutive transport to the duration and length of the same state.

Composition with Barbero--Immirzi preserves additional information:
the relative sign of the channels. Hexad--octad incidence and the
calendar determine the terms of the angular functional; the inverse
of $f$ expresses it on $\mathcal V$.
Theorem~\ref{iii:thm:barbero-factorization} proves this
factorization. Channel interchange relative to a fixed orientation
mark reverses the functional's sign; simultaneous conjugation of
operator and mark preserves it. The constitutive response thereby
transmits the information required by its angular readout.

% BEGIN SINTESIS_ADITIVA_20260922
\par
Section~\ref{delta22:iii:restitucion} assembles the spectral potential of constitutive response, global channel recovery through normalized speed and Barbero--Immirzi, and stability bounds. On the generated image, retaining labels and bases, the composition recovers the angular pair, action section, and periodic reading of the dodecaphasic record.
% END SINTESIS_ADITIVA_20260922

Section~\ref{iii:sec:incidencias-espectrales} makes the normalizations
of this composition explicit. Positive factorization separates the
common response factor from its oriented unimodular component;
together these data restore all four constitutive coordinates.
The Fourier representation then distinguishes the modes of the
$90/120$ incidences, retaining the chain's signed weights and a
declared transform normalization. This organization compares
antecedent realizations with the operator used in the article,
preserving each realization's function and domain.

Appendix~\ref{iii:hist:antecedentes} also develops the finite realization
on positions, modes, and phase: it proves equivariance and decomposition
into twenty-seven components, retains the phases of their eigenvalues,
and evaluates parametric sensitivity. The antecedent angular readouts
are reproduced with their declared operations and data. Comparison with
the constitutive readouts in the main text distinguishes the weights,
normalizations, and domains of each realization.

\subsection{Polarization, charge and dimensional realization}

The electrical development also relies on two results of the
arithmetic sheets. The mismatch between nonadic and decimal capacities
produces a uniformly bounded accumulated discrepancy; its temporal
difference defines the polarization increment. In addition, the
pairing of additive and multiplicative fibres in $D_9^3$ determines
the tritic mode and the commutator norm $\sqrt3/4$ of their
projections. Section~\ref{iii:sec:vacancias-prefactor} preserves
the proofs and domains of both constructions. Their relation to the
limit $3/4$ of the temporal response specifies the datum they share
without identifying their operator spaces.

Section~\ref{iii:sec:carga} composes impedance, $\alpha$ and
Planck's constant through the unique positive section satisfying
$Z_{\rm vac}e_a^2=2\alpha h_a$. It yields the von Klitzing
resistance, conductance quantum, magnetic flux and Josephson constant,
with exact cross-identities. Action return leaves the first two
invariant and transforms the latter two with opposite powers of the
action ratio. This electrical content follows from the preceding
sections rather than constituting a list of independent constants.

The final evaluation retains the action, charge and speed lines,
as well as the transformations of their coordinates under changes
of units. The table in Section~\ref{iii:sec:evaluacion} presents
the internal coordinates obtained by evaluation. Physical comparison is subsequently formulated
with the dimensional chart and coupling regime stated. The argument
is organized by one common question: which structure determines the
response, which equations it imposes and which information can be
recovered from the values thus obtained.
